% Ad Quintum Fratrem 3.9 (ad-quintum-fratrem-03-09) --- English translation
% Translated by Claude (Anthropic) from the Latin (Perseus).
% Style guide: STYLE.md.

\ciceroLetterOpener{MARCVS QVINTO FRATRI SALVTEM}

\ciceroSection{1} On Gabinius there was no need for me to do any of those things which you so affectionately thought up. May the earth open up for me then ---! \textit{[Greek: tote moi chanoi]} I acted, as everyone agrees, with the utmost gravity, and what I did, with the utmost lenity. I neither pressed the man down nor gave him a lift; I was a stinging witness, and for the rest kept quiet. The outcome of the trial --- foul and ruinous --- I bore as lightly as possible; and indeed this much good is now at last running over to me from it, that under these miseries of the commonwealth and this licence of the reckless, by which I used formerly to be broken, I am now not even moved. For there is nothing more abandoned than these men, than these times.

\ciceroSection{2} And so, since no further pleasure can be taken from public life, I do not know why I should fume. My letters, our studies, my leisure and the country houses delight me, and above all our boys. Milo alone vexes me. But I do wish the consulship would bring an end to the matter; in which I shall strain myself no less than I strained myself in our own, and you from where you are will help, as you are doing. As for that affair, except for what raw force may snatch from us, all is in good order; the family finances --- those I fear. \textit{[Greek: ho de mainetai ouk et' anekt\=os]} ``but he rages no longer bearably'' --- he, who is fitting up a school of gladiators at a hundred thousand sesterces! In this one matter I shall bear his rashness as I can, and that you may bear it is a charge upon your sinews.

\ciceroSection{3} As for the unsettled times of the coming year, I had not meant you to understand any domestic alarm of mine, but rather the general condition of the commonwealth; in which, even though I am attending to nothing, I can hardly avoid being anxious about anything. How cautious I would wish you to be in writing, you may gather from this --- that I am not even writing to you about those things which are openly upset in public affairs, for fear that, if my letter were intercepted, it might offend somebody's feelings. So I wish you free of domestic anxiety; in public anxiety I know how worried you tend to be. I see our Messala consul, if through an interrex, then without a trial, if through a dictator, still without danger. He has nothing they hate; Hortensius's warmth will carry much weight; Gabinius's acquittal is reckoned a law of impunity. \textit{[Greek: en parerg\=oi]} (by the way) on the dictatorship nothing as yet has been done. Pompey is away; Appius is stirring it all; Hirrus is preparing for it; many vetoers are being counted; the people do not care; the leading men do not want it; I am keeping quiet.

\ciceroSection{4} As for the slaves you promise me --- I do love you very much, and indeed I am, as you write, undermanned both at Rome and on the estates; but for heaven's sake, my brother, do not think of anything that touches my convenience unless it is at the highest convenience to yourself and the highest ease of your own means.

\ciceroSection{5} About Vatinius's letter I laughed; but I know I am so closely observed by him that those hatreds of his I not only must swallow but must also digest.

\ciceroSection{6} As to your urging me to finish it: I have finished --- to my taste at least, the thing is rather sweet --- an \textit{epos} \textit{[Greek: epos]} (epic poem) addressed to Caesar; but I am hunting for a well-set-up courier, lest there happen to it what happened to your Erigone, for whom alone, with Caesar as general, the route out of Gaul was unsafe.

\ciceroSection{7} What then? If I did not have good masonry, should I pull the building down? It is what pleases me more every day; and above all the lower colonnade and its rooms are coming out right. As for the place at Arcanum, there is need of Caesar, or, by Hercules, of someone with even better taste; for those statues of yours and the wrestling-ground and the fish-pond and the Nile are work for many a Philotimus, not for any Diphilus. But I shall go after these things myself, and shall send people, and shall give instructions.

\ciceroSection{8} About Felix's will you will complain the more, if you know the facts. For the tablets which he supposed he was sealing --- in which he held you most firmly down for a twelfth share --- those by an error both his own and his slave Sicura's he did not seal; the ones he did not mean to seal, he sealed. \textit{[Greek: all' oim\=oz\=et\=o]} ``well, let him howl for it!'' --- only let us be well. Cicero I love both as you ask and as he deserves and as I ought; I am sending him away from me, both to keep from drawing him off from his teachers, and because his mother Porcia is not going away --- without whom I am quite afraid of the boy's greediness at table. But we are very much together for all that. I have written back on every point. My sweetest and best brother, farewell.
